\begin{frame}{왜 ``직선이 무수히 많다'': 스케일 불변성}
  ``$w^\top x + b = 0$을 만족하는 $(w,b)$는 하나뿐이지 않나?''
  \begin{itemize}
    \item 답: 직선 \textbf{자체}는 $(w,b)$를 (음수가 아닌) 상수 $k$로
      곱해도 변하지 않는다:
      \[
      k\,w^\top x + k\,b = 0 \;\iff\; w^\top x + b = 0
      \]
    \item $(w,b)$와 $(3w,3b)$는 \textbf{정확히 같은 결정 경계}.
      이를 \kb{스케일 불변성}(scale invariance)이라 부른다.
  \end{itemize}
  \vskip0.6em
  \begin{block}{이 사실이 왜 중요한가}
    ``$\|w\|$를 최소화하자''를 \textbf{제약 없이} 풀면
    $w\to 0$으로 마음대로 쪼그라든다 — 경계선은 한 치도 안 움직이는데
    ``마진'' $2/\|w\|$만 끝없이 부풀러진다.
    \vskip0.3em
    즉 스케일 불변성을 그대로 두면 마진 최대화는 \textbf{정해지지 않은} 문제.
    $\|w\|$를 ``의미 있는 크기''로 만들려면 스케일을 못 박는 조건이 필요하다.
  \end{block}
\end{frame}
